
\subsubsection{Interpreting the Synthetic Data Results}

The synthetic data results demonstrate that the statistical methodology can successfully detect: (1) severe documentation gaps (7\% compliance with stringent CI gates), (2) non-uniform component deficiencies (licensing weakest despite mechanical simplicity), and (3) mechanism specificity (commits correlate while contributors and issues do not). These patterns validate the study design and analysis pipeline for deployment on real data.

\subsubsection{Limitations}

\textbf{L1: Proof-of-Concept Synthetic Data (SEVERITY: HIGH)}  
The current implementation uses synthetic data matching expected distributions to validate statistical methodology and gate logic. Results demonstrate that the detection methodology works, but the actual compliance rate (7\%) and correlation magnitude ($\rho$ = 0.951) are hypothetical until confirmed on real data. Production deployment requires HuggingFace Hub API sampling, GitHub API activity collection, and manual DCS\_3 coding by trained coders. The study design is valid; the specific numerical results require empirical confirmation.

\textbf{L2: Cross-Sectional Correlation, Not Causation (SEVERITY: MEDIUM)}  
The correlation analysis measures activity and documentation at a single timepoint (T0+90), precluding causal inference. Directionality cannot be determined: whether commits lead to better documentation (workflow integration) or documentation enables more commits (better onboarding) remains unclear. Temporal precedence (commits measured T0--T90, DCS measured at T90) provides weak directionality evidence, but correlation remains the strongest defensible claim. Causal testing requires longitudinal analysis or randomized intervention.

\textbf{L3: Single Platform Specification (HuggingFace Only) (SEVERITY: MEDIUM)}  
The study design targets HuggingFace Datasets Hub repositories and may not generalize to other platforms (Papers with Code, OpenML, Zenodo) or earlier time periods. HuggingFace is the largest public ML dataset platform, making it representative of contemporary open ML practices, but platform-specific norms may influence compliance rates. Multi-platform replication would test whether findings are HuggingFace-specific or field-wide.

\textbf{L4: 3-Component Subset (Not Full Rubric) (SEVERITY: LOW)}  
DCS\_3 measures only data context, preprocessing, and licensing (3 of 14 Rondina components). Full rubric compliance may differ if other components (ethics, intended use, known limitations) are better or worse documented. However, DCS\_3 components are foundational (Rondina factor 1: Core Documentation) and most critical for reproducibility. The 3-component subset was chosen for feasibility, not arbitrary restriction.

\subsubsection{Implications for Real-World Deployment}

If the synthetic data patterns hold in real repositories, the findings would suggest that documentation compliance is driven by workflow integration during active development rather than by framework awareness or community size. Three intervention directions follow:

\textbf{1. Commit-Triggered Documentation Prompts:} Repository platforms could implement pre-commit hooks or CI/CD checks prompting documentation updates when commits modify dataset files. This would exploit the commit-documentation correlation observed in synthetic data: repositories with active development would encounter prompts precisely when most likely to comply.

\textbf{2. Automated Licensing Templates:} If 73\% of real repositories similarly lack LICENSE files, automated license selection integrated into repository upload flow could address this barrier with minimal friction.

\textbf{3. Documentation Health Scores:} Visible documentation quality badges could leverage social pressure among active repositories, if the synthetic pattern of commit-documentation correlation holds in practice.

These interventions require empirical validation through A/B testing on real repositories. The synthetic data results provide proof-of-concept that the statistical framework can detect intervention effects if they exist.
